\documentclass[english,parskip=half-,paper=a4]{scrartcl}
\usepackage[margin=1.5cm,includeheadfoot]{geometry}
\RequirePackage[default]{lato}
\RequirePackage[math-style=french,bold-style=ISO]{lete-sans-math}
\setsansfont{Fira Sans}[Scale=MatchLowercase]
\setmonofont{CascadiaMono}[Scale=MatchLowercase]

\usepackage{tkz-interval}

\usepackage{babel}
\usepackage{enumitem}

\usepackage{minted-code}
\usepackage{fancyvrb}
\usepackage{xcolor,url}
\usepackage[colorlinks]{hyperref}
\setlength\parindent{0pt}

\def\Lcs#1{\texttt{\textbackslash#1}\index{#1@\texttt{\textbackslash#1}}}
\def\Largb#1{\texttt{\{}\textit{#1}\texttt{\}}}
\def\Lpack#1{\texttt{#1}\index{#1@\texttt{#1}}\index{Package!#1@\texttt{#1}}}
\def\Loption#1{\texttt{#1}\index{#1@\texttt{#1}}\index{Option!#1@\texttt{#1}}}

\title{Interval brackets, with Ti\textit{k}Z (1.0a)}
\date{October 12, 2025}
\author{Cédric Pierquet\thanks{\protect\url{cpierquet@outlook.fr}}}
\raggedbottom

\def\IntervtestA{-4 ; \dfrac{1}{7}}
\def\IntervtestB{\dfrac{1}{2} ; \sqrt{7}}

\begin{document}

\maketitle

\tableofcontents

\section{Introduction}

\subsection{Loading}

To use \Lpack{tkz-interval}, simply use:

\begin{codeblockA}[title=Loading]
\usepackage{tkz-interval}
\end{codeblockA}

Used packages are \Lpack{tikz} and \Lpack{simplekv}.

\subsection{Idea(s)}

The idea behind this package is to provide a command for representing an interval with brackets created using Ti\textit{k}Z, with automatic size management and customization of thickness and depth, for examples.

\smallskip

It's also possible to specify around spaces and optional \textit{overlap} for open bracket.

\smallskip

This documentation was created using the \Lpack{minted-code} package.

\subsection{Limitations}

Spaces around and inside seems to be fine, but not perfect and not exactly like \textit{normal} version\ldots

\subsection{First example(s)}

\begin{codeblockA}[title=Samples intervals]
\def\IntervtestA{ -4 ; \dfrac{1}{7}       }
\def\IntervtestB{ \dfrac{1}{2} ; \sqrt{7} }
\end{codeblockA}

\begin{exampleM}{}
%'normal' version
$x \in \left[ \IntervtestA \right] \cup \left[ \IntervtestB \right[$.\\
$x \in \left] \IntervtestA \right[ \cup \left[ \IntervtestB \right[$.\\
$x \in \left[ \IntervtestA \right[ \cup \left[ \IntervtestB \right[$.\\
$x \in \left] \IntervtestA \right] \cup \left[ \IntervtestB \right[$.\\
\end{exampleM}

\begin{exampleM}{}
%with tkz-interval package
\setKVdefault[tkzinterv]{left color=teal,right color=red}
$x \in \tkzinterval{\IntervtestA} \cup \tkzinterval[open]{\IntervtestB}$.\\
$x \in \tkzinterval[open]{\IntervtestA} \cup \tkzinterval[open]{\IntervtestB}$.\\
$x \in \tkzinterval[open right]{\IntervtestA} \cup \tkzinterval[open]{\IntervtestB}$.\\
$x \in \tkzinterval[open left]{\IntervtestA} \cup \tkzinterval[open]{\IntervtestB}$.
\end{exampleM}

\subsection{History}

\texttt{1.0a}: Initial version

\pagebreak

\section{Usage}

\subsection{Global lengths, local keys}

\begin{codeblockA}[title=Global lengths]
\setlength\tkzintervhoriz{length}   %brackets width        (5pt)
\setlength\tkzintervwidth{length}   %brackets thickness  (0.7pt)
\end{codeblockA}

\begin{codeblockA}[title=Local keys]
\setKVdefault[tkzinterv]{key1 = new default value,...}
\end{codeblockA}

\subsection{Main macro}

\begin{codeblockA}[title=Macro (english usage)]
\tkzinterval[type,keys]{lower limit ; upper limit}
%type can be open / open left / open right (closed by default)
\end{codeblockA}

\begin{codeblockA}[title=Macro (french usage)]
\tkzintervalle[type,clés]{borne inf ; borne sup}
%le type peut-être FF (défaut) / OO / OF / FO
\end{codeblockA}

\subsection{Keys}

\begin{itemize}[noitemsep]
	\item \texttt{vsep}: vertical space (in \textit{unit}) between box and brackets \hfill(\texttt{0.1em})
	\item \texttt{overlap}: 'overlap 'effect' (\texttt{lr-l-r-none}) for open bracket \hfill(\texttt{lr})
	\item \texttt{overlap factor}: overlap factor(s) (\texttt{lr-l/r}) for open brackets \hfill(\texttt{0.5})
	\item \texttt{hsep}: horizontal margin, as a percentage of the width of the brackets\hfill(\texttt{0.75})
	\item \texttt{space before}: space before interval \hfill(\texttt{\textbackslash,})
	\item \texttt{space after}: space after interval \hfill(\texttt{\textbackslash,})
	\item \texttt{nobeforeafter}: suppress space around interval \hfill(\texttt{false})
\end{itemize}

\section{Influence of some parameters}

\subsection{Default values}

\begin{codeblockA}[title=Default values]
\setlength\tkzintervhoriz{5pt}
\setlength\tkzintervwidth{0.7pt}
\setKVdefault[tkzinterv]{hsep=0.75}
\end{codeblockA}

\setlength\tkzintervhoriz{5pt}
\setlength\tkzintervwidth{0.7pt}
\setKVdefault[tkzinterv]{hsep=0.75}

\begin{exampleM}{}
Let's consider the interval $I=\tkzinterval[open,overlap=lr]{\IntervtestA}$.\par
Let's consider the interval $I=\tkzinterval[open,overlap=l]{\IntervtestA}$.\par
Let's consider the interval $I=\tkzinterval[open,overlap=r]{\IntervtestA}$.\par
Let's consider the interval $I=\tkzinterval[open,overlap=none]{\IntervtestA}$.
\end{exampleM}

\begin{exampleM}{}
Let's consider the interval $I=\tkzinterval{ \IntervtestA }$.\\
Let's consider the interval $I=\left[ \IntervtestA \right]$.
\end{exampleM}

\begin{exampleM}{}
{\huge Let's consider the interval $\mathcolor{red}{I=\tkzinterval{\IntervtestA}}$.}
\end{exampleM}

\begin{exampleM}{}
Let's consider the interval $I=\tkzinterval[open left,overlap=l]{\IntervtestA}$.\par
Let's consider the interval $I=\tkzinterval[open left,overlap=none]{\IntervtestA}$.
\end{exampleM}

\begin{exampleM}{}
Let's consider the interval $I=\tkzinterval[open right,overlap=r]{\IntervtestA}$.\par
Let's consider the interval $I=\tkzinterval[open right,overlap=none]{\IntervtestA}$.
\end{exampleM}

\subsection{Custom values (v1)}

\begin{codeblockA}[title=Custom values]
\setlength\tkzintervhoriz{6pt}
\setlength\tkzintervwidth{0.95pt}
\setKVdefault[tkzinterv]{hsep=0.5,vsep=1mm}
\end{codeblockA}

\setlength\tkzintervhoriz{6pt}
\setlength\tkzintervwidth{0.95pt}
\setKVdefault[tkzinterv]{hsep=0.5,vsep=1mm}

\begin{exampleM}{}
Let's consider the interval $I=\tkzinterval[open,overlap=lr]{\IntervtestA}$.\par
Let's consider the interval $I=\tkzinterval[open,overlap=l]{\IntervtestA}$.\par
Let's consider the interval $I=\tkzinterval[open,overlap=r]{\IntervtestA}$.\par
Let's consider the interval $I=\tkzinterval[open,overlap=none]{\IntervtestA}$.
\end{exampleM}

\begin{exampleM}{}
Let's consider the interval $I=\tkzinterval{ \IntervtestA }$.\\
Let's consider the interval $I=\left[ \IntervtestA \right]$.
\end{exampleM}

\begin{exampleM}{}
{\huge Let's consider the interval $\mathcolor{red}{I=\tkzinterval{\IntervtestA}}$.}
\end{exampleM}

\begin{exampleM}{}
Let's consider the interval $I=\tkzinterval[open left,overlap=l]{\IntervtestA}$.\par
Let's consider the interval $I=\tkzinterval[open left,overlap=none]{\IntervtestA}$.
\end{exampleM}

\begin{exampleM}{}
Let's consider the interval $I=\tkzinterval[open right,overlap=r]{\IntervtestA}$.\par
Let's consider the interval $I=\tkzinterval[open right,overlap=none]{\IntervtestA}$.
\end{exampleM}

\subsection{Custom values (v2)}

\begin{codeblockA}[title=Custom values]
\setlength\tkzintervhoriz{1.75mm}
\setlength\tkzintervwidth{0.65mm}
\setKVdefault[tkzinterv]{hsep=1,vsep=0.125em}
\end{codeblockA}

\setlength\tkzintervhoriz{1.75mm}
\setlength\tkzintervwidth{0.65mm}
\setKVdefault[tkzinterv]{hsep=1,vsep=0.125em}

\begin{exampleM}{}
Let's consider the interval $I=\tkzinterval[open,overlap=lr]{\IntervtestA}$.\par
Let's consider the interval $I=\tkzinterval[open,overlap=l]{\IntervtestA}$.\par
Let's consider the interval $I=\tkzinterval[open,overlap=r]{\IntervtestA}$.\par
Let's consider the interval $I=\tkzinterval[open,overlap=none]{\IntervtestA}$.
\end{exampleM}

\begin{exampleM}{}
Let's consider the interval $I=\tkzinterval{ \IntervtestA }$.\\
Let's consider the interval $I=\left[ \IntervtestA \right]$.
\end{exampleM}

\begin{exampleM}{}
{\huge Let's consider the interval $\mathcolor{red}{I=\tkzinterval{\IntervtestA}}$.}
\end{exampleM}

\begin{exampleM}{}
Let's consider the interval $I=\tkzinterval[open left,overlap=l]{\IntervtestA}$.\par
Let's consider the interval $I=\tkzinterval[open left,overlap=none]{\IntervtestA}$.
\end{exampleM}

\begin{exampleM}{}
Let's consider the interval $I=\tkzinterval[open right,overlap=r]{\IntervtestA}$.\par
Let's consider the interval $I=\tkzinterval[open right,overlap=none]{\IntervtestA}$.
\end{exampleM}

\subsection{Overlap factor}

\begin{codeblockA}[title=Custom values]
\setlength\tkzintervhoriz{5pt}
\setlength\tkzintervwidth{0.9pt}
\end{codeblockA}

\setlength\tkzintervhoriz{5pt}
\setlength\tkzintervwidth{0.9pt}

\begin{exampleM}{}
\setlength\fboxsep{0.5pt}
.\fbox{$\tkzinterval[open,overlap factor=0.25] {\dfrac{1}{2};\sqrt{7}}$}.\\
.\fbox{$\tkzinterval[open,overlap factor=0.5]  {\dfrac{1}{2};\sqrt{7}}$}.\\
.\fbox{$\tkzinterval[open,overlap factor=0.75] {\dfrac{1}{2};\sqrt{7}}$}.\\
.\fbox{$\tkzinterval[open,overlap factor=1]    {\dfrac{1}{2};\sqrt{7}}$}.\\
.\fbox{$\tkzinterval[open,overlap factor=0/1]  {\dfrac{1}{2};\sqrt{7}}$}.
\end{exampleM}

\subsection{Samples (v2)}

\begin{codeblockA}[title=Custom values]
\setlength\tkzintervhoriz{5pt}
\setlength\tkzintervwidth{0.9pt}
\setKVdefault[tkzinterv]{hsep=0.75,left color=teal,right color=red}
\end{codeblockA}

\setlength\tkzintervhoriz{5pt}
\setlength\tkzintervwidth{0.9pt}
\setKVdefault[tkzinterv]{hsep=0.75,left color=teal,right color=red}

\begin{exampleM}{}
\setKVdefault[tkzinterv]{hsep=0.75,left color=teal,right color=red}
$x \in \tkzinterval[]{4;\dfrac{100}{7}} \cup \tkzinterval[open]{\dfrac{1}{2};\sqrt{7}}$.\hfill
$x \in \tkzinterval[open]{4;\dfrac{100}{7}} \cup \tkzinterval[open]{\dfrac{1}{2};\sqrt{7}}$.\hfill\null

\hfill$x \in \tkzinterval[open right]{4;\dfrac{100}{7}} \cup \tkzinterval[open]{\dfrac{1}{2};\sqrt{7}}$.\hfill
$x \in \tkzinterval[open left]{4;\dfrac{100}{7}} \cup \tkzinterval[open]{\dfrac{1}{2};\sqrt{7}}$.
\end{exampleM}

\begin{exampleM}{}
$x \in \tkzinterval[]{4;\dfrac{100}{7}} \cup \tkzinterval[]{\dfrac{1}{2};\sqrt{7}}$.\hfill
$x \in \tkzinterval[open]{4;\dfrac{100}{7}} \cup \tkzinterval[]{\dfrac{1}{2};\sqrt{7}}$.\hfill\null

\hfill$x \in \tkzinterval[open right]{4;\dfrac{100}{7}} \cup \tkzinterval[]{\dfrac{1}{2};\sqrt{7}}$.\hfill
$x \in \tkzinterval[open left]{4;\dfrac{100}{7}} \cup \tkzinterval[]{\dfrac{1}{2};\sqrt{7}}$.
\end{exampleM}

\begin{exampleM}{}
$x \in \tkzinterval[]{4;\dfrac{100}{7}} \cup \tkzinterval[open left]{\dfrac{1}{2};\sqrt{7}}$.\hfill
$x \in \tkzinterval[open]{4;\dfrac{100}{7}} \cup \tkzinterval[open left]{\dfrac{1}{2};\sqrt{7}}$.\hfill\null

\hfill$x \in \tkzinterval[open right]{4;\dfrac{100}{7}} \cup \tkzinterval[open left]{\dfrac{1}{2};\sqrt{7}}$.\hfill
$x \in \tkzinterval[open left]{4;\dfrac{100}{7}} \cup \tkzinterval[open left]{\dfrac{1}{2};\sqrt{7}}$.
\end{exampleM}

\begin{exampleM}{}
$x \in \tkzinterval[]{4;\dfrac{100}{7}} \cup \tkzinterval[open right]{\dfrac{1}{2};\sqrt{7}}$.\hfill
$x \in \tkzinterval[open]{4;\dfrac{100}{7}} \cup \tkzinterval[open right]{\dfrac{1}{2};\sqrt{7}}$.\hfill\null

\hfill$x \in \tkzinterval[open right]{4;\dfrac{100}{7}} \cup \tkzinterval[open right]{\dfrac{1}{2};\sqrt{7}}$.\hfill
$x \in \tkzinterval[open left]{4;\dfrac{100}{7}} \cup \tkzinterval[open right]{\dfrac{1}{2};\sqrt{7}}$.
\end{exampleM}

\end{document}

